\begin{abstract}
Engineering education research identifies estimation as a critical skill reliant on expert cognitive processes like mental simulation, yet it remains under-taught. The Mettle (Modelling-based Estimation Learning Environment) platform scaffolds this learning but lacks a mechanism for providing dynamic, individualized Socratic feedback. This study addresses this gap by selecting and rigorously benchmarking a Large Language Model (LLM) to serve as an integrated, adaptive Socratic chatbot within Mettle. A multi-dimensional evaluation framework was implemented to assess candidate models along axes of pedagogical quality, contextual adaptability, cost-efficiency, and technical performance. Five models—proprietary systems (OpenAI gpt-4o-mini, Anthropic claude-3-sonnet-20240229, Google gemini-2.5-flash, Cohere command-r-08-2024) and an open-source model (Meta Llama-4-Maverick-17B-128E-Instruct)—were evaluated. Using a purpose-built, reproducible Python benchmarking harness and a curated 50-prompt bank, we employed a mixed-methods approach combining automated performance metrics with human-rated rubrics to score Socratic dialogue quality. The analysis yields a comparative ranking, highlighting the quantitative trade-offs between pedagogical effectiveness and operational cost. The findings provide a single, data-driven recommendation for the optimal LLM to enhance Mettle, offering a scalable pathway to personalized learning support that encourages the mental simulation essential to expert engineering practice.
\end{abstract}

\begin{keywords}
LLM Benchmarking, Socratic Tutoring, Engineering Education, Model Evaluation, Mettle, Reproducible Research
\end{keywords}